\documentclass[12pt,a4paper]{article}
\usepackage[utf8]{inputenc}
\usepackage[T1]{fontenc}
\usepackage[spanish,es-noshorthands]{babel}
\usepackage{amsmath,amssymb,geometry,fancyhdr,xcolor,listings,tikz,array}
\geometry{margin=2.4cm,headheight=16pt}
\pagestyle{fancy}\fancyhf{}\lhead{UPCT}\rhead{Cálculo diferencial}\lfoot{Prácticas de wxMaxima}\rfoot{\thepage}
\definecolor{azul}{RGB}{23,65,105}
\lstset{basicstyle=\ttfamily\small,backgroundcolor=\color{black!4},frame=single,rulecolor=\color{black!15},breaklines=true,columns=fullflexible}

\begin{document}
\begin{center}{\large Práctica de laboratorio II}\\[5pt]{\LARGE\bfseries\color{azul} Derivadas, diferenciabilidad y gradiente}\end{center}
\noindent Nombre y apellidos: \hrulefill\par\medskip
\noindent\textbf{Objetivos.} Distinguir derivadas según vectores y direcciones, calcular parciales, comprobar diferenciabilidad y comprender la diferencial y el gradiente. Ejecutar las celdas con Mayús+Intro. Justificar cada conclusión matemáticamente.
\begin{lstlisting}
kill(all)$
\end{lstlisting}
\section{Derivadas según un vector}
Sea $f:U\subset\mathbb R^2\to\mathbb R$, con $U$ abierto, $P=(a,b)$ y $w=(u,v)\ne0$. Definimos
\[D_wf(P)=\lim_{t\to0}\frac{f(P+tw)-f(P)}{t}.\]
No es necesario que $w$ sea unitario. Para $c\ne0$, si existe la derivada, $D_{cw}f(P)=cD_wf(P)$ (cambio de parámetro). Esto no implica aditividad respecto del vector.
\subsection*{Ejemplo guiado} Sea $f(x,y)=x^2+xy+2y^2$, $P=(1,2)$ y $w=(3,4)$.
\begin{lstlisting}
f(x,y):=x^2+x*y+2*y^2$
q:expand((f(1+3*t,2+4*t)-f(1,2))/t)$
limit(ratsimp(q),t,0);
\end{lstlisting}
El resultado es $\boxed{48}$: las parciales en $P$ son $4$ y $9$. Repetir con $-w$ y $2w$.
\section{Derivadas según una dirección}
Llamaremos derivada según una dirección orientada a la derivada según el vector unitario $e=w/\|w\|$. En particular,
\[\boxed{D_ef(P)=\frac{D_wf(P)}{\|w\|}}.\]
La normalización permite medir la variación por unidad de distancia recorrida. En el ejemplo anterior $\|w\|=5$.
\begin{lstlisting}
norma:sqrt(3^2+4^2)$
limit(ratsimp((f(1+3*t/norma,2+4*t/norma)-f(1,2))/t),t,0);
\end{lstlisting}
Resultado: $48/5$. Comparar con la dirección opuesta. Para una dirección de ángulo $\theta$, usar $e=(\cos\theta,\sin\theta)$.
\section{Derivadas parciales}
Son las derivadas según los vectores de la base canónica:
\[f_x(a,b)=D_{(1,0)}f(a,b),\qquad f_y(a,b)=D_{(0,1)}f(a,b).\]
Se pueden calcular por definición o mediante derivación simbólica, manteniendo fija la otra variable.
\begin{lstlisting}
limit(ratsimp((f(1+t,2)-f(1,2))/t),t,0);
limit(ratsimp((f(1,2+t)-f(1,2))/t),t,0);
fx:diff(f(x,y),x);
fy:diff(f(x,y),y);
subst([x=1,y=2],[fx,fy]);
\end{lstlisting}
Se obtienen $4$ y $9$. Actividad: comparar la derivación simbólica con los cocientes incrementales en otro punto.
\section{Diferenciabilidad en un punto}
La función es diferenciable en $P$ si existe una aplicación lineal $L:\mathbb R^2\to\mathbb R$ tal que
\[f(P+h)-f(P)=L(h)+r(h),\qquad \lim_{h\to0}\frac{r(h)}{\|h\|}=0.\]
Necesariamente $L(h)=f_x(P)h_1+f_y(P)h_2$. Esta aplicación se denomina diferencial $df_P$. La diferenciabilidad implica continuidad y existencia de todas las derivadas según vectores, con $D_wf(P)=df_P(w)$.
Para nuestra función, el resto es cuadrático.
\begin{lstlisting}
R:expand(f(1+h,2+k)-f(1,2)-4*h-9*k);
Rpolar:expand(subst([h=r*cos(th),k=r*sin(th)],R));
limit(ratsimp(Rpolar/r),r,0,plus);
\end{lstlisting}
Como $|h^2+hk+2k^2|\le(5/2)(h^2+k^2)$, el cociente del resto por la norma tiende uniformemente a cero. Un límite polar para cada ángulo fijo no bastaría por sí solo.
\section{Condición suficiente de diferenciabilidad}
Si las derivadas parciales existen en un entorno abierto de $P$ y son continuas en $P$, entonces $f$ es diferenciable en $P$. Es una condición suficiente, no necesaria. El polinomio anterior cumple el criterio.
\subsection*{Contraejemplo} Definimos $g(0,0)=0$ y, fuera del origen,
\[g(x,y)=\frac{x^3}{x^2+y^2}.\]
Es continua porque $|g(x,y)|\le |x|$. Todas las derivadas según vectores existen, pero no forman una aplicación lineal.
\begin{lstlisting}
g(x,y):=x^3/(x^2+y^2)$
limit(ratsimp(g(t,0)/t),t,0);
limit(ratsimp(g(0,t)/t),t,0);
limit(ratsimp(g(t,t)/t),t,0);
limit(ratsimp((g(t,t)-t)/(sqrt(2)*abs(t))),t,0,plus);
gx:ratsimp(diff(g(x,y),x));
limit(subst(x=0,gx),y,0);
\end{lstlisting}
Las parciales son $1,0$; la derivada según $(1,1)$ es $1/2$, no $1$. El cociente del resto por la norma sobre $t>0$ vale $-1/(2\sqrt2)$. Además, $g_x(0,y)=0$ para $y\ne0$, así que $g_x$ no es continua en el origen. En los comandos, $g$ representa solo la fórmula fuera del origen.
\section{Gradiente, diferencial e interpretación geométrica}
Si $f$ es diferenciable, su gradiente es $\nabla f(P)=(f_x(P),f_y(P))$ y
\[\boxed{df_P(h)=\nabla f(P)\cdot h,\quad D_ef(P)=\nabla f(P)\cdot e}.\]
Para $\|e\|=1$, el máximo es $\|\nabla f(P)\|$, alcanzado en la dirección del gradiente si no es nulo; el mínimo es su opuesto. El gradiente es perpendicular a la curva de nivel regular. El plano tangente a la gráfica es $z=f(P)+df_P(x-a,y-b)$.
En el ejemplo $f(1,2)=11$, $\nabla f(1,2)=(4,9)$.
\begin{lstlisting}
grad:[4,9]$
modulo:sqrt(4^2+9^2);
emax:grad/modulo;
plano:11+4*(x-1)+9*(y-2);
wxplot3d([f(x,y),plano],[x,0,2],[y,1,3])$
wxplot2d(4*cos(th)+9*sin(th),[th,0,2*%pi])$
\end{lstlisting}
El máximo es $\sqrt{97}$ y el mínimo $-\sqrt{97}$. Una dirección tangente a la curva de nivel es $(-9,4)/\sqrt{97}$, con derivada cero.
\section{Trabajo individual}
En cada ejercicio, utilizar wxMaxima para verificar los cálculos y escribir una justificación independiente cuando se estudie diferenciabilidad.
\begin{enumerate}
\item $f=x^2+3xy-y^2$, $P=(1,-1)$, $w=(2,-1)$. Calcular parciales, derivada según el vector y según su dirección, gradiente y diferencial.
\item $f=e^{xy}$, $P=(1,0)$, $w=(1,2)$. Justificar diferenciabilidad y hallar el plano tangente.
\item $f(0,0)=0$, $f(x,y)=xy/(x^2+y^2)$ fuera del origen. Calcular las parciales y estudiar continuidad y diferenciabilidad.
\item $f(0,0)=0$, $f(x,y)=x^2y/(x^2+y^2)$ fuera del origen. Estudiar todas las derivadas según vectores y la diferenciabilidad.
\item $f(0,0)=0$, $f(x,y)=x^2\sin(1/x)$ si $x\ne0$ y $0$ si $x=0$. Probar diferenciabilidad en el origen pese a la discontinuidad de $f_x$.
\item $f=x^2+4y^2$, $P=(2,1)$. Hallar las direcciones de máximo, mínimo y variación nula; representar la gráfica y su plano tangente.
\end{enumerate}
\section{Cuestiones de comprensión}
\begin{enumerate}
\item ¿Qué cambia al multiplicar un vector por un escalar negativo? ¿Y al normalizarlo?
\item ¿La existencia de parciales implica continuidad? ¿La existencia de todas las derivadas según vectores implica diferenciabilidad?
\item ¿Por qué la continuidad de las parciales es suficiente pero no necesaria?
\item ¿Qué relación existe entre gradiente, curva de nivel y plano tangente?
\end{enumerate}
\noindent\textbf{Entrega.} Guardar el cuaderno ejecutado con nombre y apellidos, incluyendo las conclusiones.
\end{document}
